\subsection{The instructions, verbatim}
\label{app:prompts}

Every model saw one of these, followed by the study image and the four options.
The sweep in Figure~\ref{fig:landmarks} runs the first six on Qwen2.5-VL-32B;
every other run in the paper uses \texttt{cot\_anyview}. Lines are wrapped to fit the page.

\paragraph{\texttt{cot\_anyview}} chain of thought, viewpoint unstated.

\begin{footnotesize}
\begin{verbatim}
You are taking the Four Mountains Test of spatial allocentric perception.

Image 1 is the STUDY view of a landscape with four mountain peaks.
The subsequent 4 images are OPTION 1 to OPTION 4.
Exactly ONE option shows the EXACT SAME mountain landscape (the same four peaks in
  the same relative spatial arrangement). It may be viewed from the same direction
  as the study image or from a different one, and the lighting/weather may differ.
The other options show different mountain landscapes.

In 2-4 sentences, compare the 3D spatial layout of the peaks (e.g. relative
  positions such as in front, behind, left, right) between the study view and the
  options, accounting for any camera rotation. Avoid lengthy itemized lists.

State your final decision on the last line as:
Final Answer: Option X
\end{verbatim}
\end{footnotesize}

\paragraph{\texttt{cot}} chain of thought, viewpoint asserted.

\begin{footnotesize}
\begin{verbatim}
You are taking the Four Mountains Test of spatial allocentric perception.

Image 1 is the STUDY view of a landscape with four mountain peaks.
The subsequent 4 images are OPTION 1 to OPTION 4.
Exactly ONE option shows the EXACT SAME mountain landscape (the same four peaks in
  the same relative spatial arrangement), simply viewed from a different viewpoint
  and under different lighting/weather.
The other options show different mountain landscapes.

In 2-4 sentences, compare the 3D spatial layout of the peaks (e.g. relative
  positions such as in front, behind, left, right) between the study view and the
  options, accounting for camera rotation. Avoid lengthy itemized lists.

State your final decision on the last line as:
Final Answer: Option X
\end{verbatim}
\end{footnotesize}

\paragraph{\texttt{mental\_rotation}} instructed to mentally rotate.

\begin{footnotesize}
\begin{verbatim}
You are taking the Four Mountains Test of allocentric spatial perception.

Image 1 is the STUDY view of a landscape with four distinct mountain peaks.
The subsequent 4 images are OPTION 1 to OPTION 4.
Exactly ONE option shows the EXACT SAME mountain landscape (the same four peaks in
  the same relative geometric layout), viewed from a shifted camera viewpoint and
  under different weather/lighting.
The other options are distractors where the relative 3D spatial arrangement of the
  peaks has been altered.

To determine the matching scene, perform mental rotation:
1. Estimate the camera perspective shift (rotation angle) between the study image
  and candidate options.
2. Mentally rotate the four peaks to verify if topological handedness (e.g.
  clockwise/counterclockwise order of peaks, which peak is opposite or between
  others) matches the study scene.
3. Ignore superficial differences in sunlight, shadows, fog, and seasonal color.

In 2-4 sentences, explain your mental rotation reasoning, then conclude on the last
  line:
Final Answer: Option X
\end{verbatim}
\end{footnotesize}

\paragraph{\texttt{anchor}} instructed to pick an anchor landmark.

\begin{footnotesize}
\begin{verbatim}
You are taking the Four Mountains Test of allocentric spatial perception.

Image 1 is the STUDY view of a landscape with four distinct mountain peaks.
The subsequent 4 images are OPTION 1 to OPTION 4.
Exactly ONE option shows the EXACT SAME mountain landscape under a camera viewpoint
  rotation and weather change.
The remaining options are distractors.

Spatial Strategy:
1. Select the single most prominent or unique landmark peak as an anchor (origin).
2. Trace the relative bearings and distances of the other 3 peaks surrounding this
  anchor.
3. Identify which option preserves this exact 3D spatial configuration around the
  anchor under the new camera angle.

In 2-4 concise sentences, explain your reasoning and conclude on the last line:
Final Answer: Option X
\end{verbatim}
\end{footnotesize}

\paragraph{\texttt{birdseye}} instructed to imagine a plan view.

\begin{footnotesize}
\begin{verbatim}
You are taking the Four Mountains Test of allocentric spatial perception.

Image 1 is the STUDY view of a landscape with four mountain peaks.
The subsequent 4 images are OPTION 1 to OPTION 4.
Exactly ONE option shows the EXACT SAME mountain landscape viewed from a different
  camera viewpoint and lighting.
The other options are distractors with altered 3D mountain configurations.

Top-Down Cognitive Mapping Strategy:
1. Imagine looking down at the four peaks from directly above (a 2D bird's-eye map).
  Note their relative positions (which forms a triangle, which is isolated, which is
  tallest).
2. For each candidate option, determine where the camera would be standing on that
  same bird's-eye map.
3. Verify which option is geometrically consistent with the study scene's top-down
  layout under the new camera angle.

In 2-4 sentences, describe the bird's-eye spatial layout and conclude on the last
  line:
Final Answer: Option X
\end{verbatim}
\end{footnotesize}

\paragraph{\texttt{elimination}} instructed to eliminate alternatives.

\begin{footnotesize}
\begin{verbatim}
You are taking the Four Mountains Test of allocentric spatial perception.

Image 1 is the STUDY view of a landscape with four mountain peaks.
The subsequent 4 images are OPTION 1 to OPTION 4.
Exactly ONE option shows the EXACT SAME mountain landscape viewed from a different
  camera angle and lighting.
The other options are geometric distractors.

Falsification Strategy:
1. Inspect each candidate option one by one to find geometric contradictions with
  the study scene (e.g. impossible relative peak heights, wrong peak ordering, or
  missing ridges).
2. Eliminate the distractor options that cannot possibly match the study landscape
  under any viewpoint rotation.
3. Select the remaining single candidate that has no geometric contradictions.

Briefly eliminate the distractors and conclude on the last line:
Final Answer: Option X
\end{verbatim}
\end{footnotesize}

\paragraph{\texttt{neutral\_anyview}} the human arm, and the only wording a person saw.

\begin{footnotesize}
\begin{verbatim}
You will see a STUDY image of a place, then 4 options.

The place contains several landmarks. Exactly ONE option shows the SAME place as the
  study image. It may be photographed from the same direction as the study image or
  from a different one, and the lighting may differ.
The other options show different places, each containing the same landmarks arranged
  differently, photographed from the same direction as the correct option.

Which option shows the same place as the study image?
Answer on the last line as:
Final Answer: Option X
\end{verbatim}
\end{footnotesize}

